\documentclass[10pt]{article}
\usepackage[T1]{fontenc}
\usepackage[utf8]{inputenc}
\usepackage[french]{babel}
\usepackage{amsmath,amssymb,mathrsfs,mathpazo}
\usepackage{geometry}
\geometry{a4paper,margin=2cm}
\usepackage{graphicx,array,tabularx,tabularray,colortbl,multirow,multicol,booktabs,enumitem,pifont,ulem,fancybox,diagbox,cancel}
\usepackage{pstricks,pstricks-add,pst-plot,pst-tree,pst-node,pst-text,pst-func,pst-math,pst-bspline,pst-3dplot,pst-eucl,pst-solides3d}
\let\pscalculate\undefined
\usepackage{pst-bezier}
\usepackage{tikz,pgfplots,tkz-tab}
\usetikzlibrary{arrows,arrows.meta,patterns,calc,positioning,trees,decorations.markings,shapes.geometric,shapes.misc,shapes.arrows,intersections,decorations.pathmorphing,decorations.pathreplacing,backgrounds,shadows,babel}
\pgfplotsset{compat=1.18}
\usepgfplotslibrary{fillbetween}
\usepackage[np]{numprint}
\usepackage{listings,makecell,pas-tableur,twoopt}
\let\barre\undefined
\usepackage[pstricks]{tablvar}
\definecolor{darkgreen}{RGB}{0,100,0}
\providecommand{\red}{\color{red}}
\providecommand{\blue}{\color{blue}}
\providecommand{\green}{\color{green}}
\providecommand{\black}{\color{black}}
\providecommand{\white}{\color{white}}
\providecommand{\gray}{\color{gray}}
\providecommand{\cyan}{\color{cyan}}
\providecommand{\magenta}{\color{magenta}}
\providecommand{\yellow}{\color{yellow}}
\providecommand{\violet}{\color{violet}}
\providecommand{\orange}{\color{orange}}
\providecommand{\euro}{\texteuro}
\providecommand{\EUR}{\texteuro}
\providecommand{\ieme}{\textsuperscript{e}}
\providecommand{\degres}{\ensuremath{^\circ}}
\providecommand{\degree}{\ensuremath{^\circ}}
\providecommand{\vv}[1]{\overrightarrow{#1}}
\providecommand{\Coord}[1]{#1}
\providecommand{\decosix}{}
\providecommand{\starredbullet}{\ensuremath{\bullet}}
\pagestyle{empty}
\setlength{\parindent}{0pt}
\usepackage{adjustbox}\DeclareUnicodeCharacter{2236}{:}\DeclareUnicodeCharacter{2190}{\ensuremath{\leftarrow}}
\providecommand{\R}{\mathbb{R}}
\providecommand{\N}{\mathbb{N}}
\providecommand{\Z}{\mathbb{Z}}
\providecommand{\Q}{\mathbb{Q}}
\providecommand{\D}{\mathbb{D}}
\providecommand{\C}{\mathbb{C}}
\providecommand{\e}{{\mathrm{e}}}
\providecommand{\vect}[1]{\overrightarrow{#1}}
\providecommand{\vectt}[1]{\overrightarrow{\mathrm{#1}}}
\providecommand{\vvt}[1]{\overrightarrow{\mathrm{#1}}}
\providecommand{\np}[2][]{{#2~#1}}
\providecommand{\vv}[1]{{\overrightarrow{#1}}}
\renewcommand{\d}{{\mathrm{d}}}
\providecommand{\degres}{{^\circ}}
\providecommand{\degre}{{^\circ}}
\providecommand{\blue}{\color{blue}}
\providecommand{\red}{\color{red}}
\providecommand{\enskip}{\quad}
\providecommand{\llbracket}{\lbrack\!\lbrack}
\providecommand{\rrbracket}{\rbrack\!\rbrack}
\providecommand{\euro}{{\text{€}}}
\providecommand{\no}{{\text{n°}}}
\providecommand{\interval}[2]{{[#1;#2]}}
\providecommand{\intervFO}[2]{{[#1;#2[}}
\providecommand{\dotfill}{\ldots}

\providecommand{\ds}{\displaystyle}
\providecommand{\textding}[1]{#1}
\providecommand{\Oij}{\(\mathrm{O};\vec{\imath},\vec{\jmath}\)}
\providecommand{\Oijk}{\(\mathrm{O};\vec{\imath},\vec{\jmath},\vec{k}\)}
\providecommand{\Ouv}{\(\mathrm{O};\vec{u},\vec{v}\)}
\providecommand{\intervFF}[2]{\ensuremath{[#1;#2]}}
\renewcommand{\euro}{\ifmmode\text{\texteuro}\else\texteuro\fi}
\DeclareUnicodeCharacter{00B0}{\ifmmode^{\circ}\else\textdegree\fi}
\DeclareUnicodeCharacter{20AC}{\ifmmode\text{\texteuro}\else\texteuro\fi}
% Énoncés : Éducation nationale. Transcription et corrigés : APMEP, Denis Vergès et contributeurs. Découpage et classement : CoopMaths, Sébastien Lozano, Éric Elter et contributeurs. Adaptation de la présentation : Maths Explorer.
% CC BY-SA : https://coopmaths.fr/www/about/
% https://coopmaths.fr/alea/static/bac/2021/tex/bac_2021_09_metropole_1.tex
\begin{document}
\noindent\textit{APMEP / CoopMaths. CC BY-SA. Présentation adaptée par Maths Explorer.}\par\medskip
\textbf{Commun à tous les candidats}

\medskip

\emph{Cet exercice est un questionnaire à choix multiples. Pour chacune des questions suivantes, une seule des quatre réponses proposées est exacte. Une réponse exacte rapporte un point. Une réponse fausse, une réponse multiple ou l'absence de réponse à une question ne rapporte ni n'enlève de point. \\Pour répondre, indiquer sur la copie le numéro de la question et la lettre de la réponse choisie.\\Aucune
justification n'est demandée.}

\medskip

\parbox{0.45\linewidth}{Le graphique ci-contre donne la représentation graphique $\mathcal{C}_f$ dans un repère orthogonal d'une fonction $f$ définie et dérivable sur $\R$.

On notera $f'$ la fonction dérivée de $f$.

On donne les points A de coordonnées (0~;~5) et B de coordonnées (1~;~20). Le point C est le point de la courbe $\mathcal{C}_f$ ayant pour abscisse $-2,5$.
La droite (AB) est la tangente à la courbe $\mathcal{C}_f$ au point A.

\smallskip

Les questions 1 à 3 se rapportent à cette même fonction $f$.}\hfill
\parbox{0.53\linewidth}{\psset{xunit=1.2cm,yunit=0.25cm}
\begin{pspicture*}(-5,-4)(1.5,22)
\psgrid[gridlabels=0pt,subgriddiv=1,gridwidth=0.05pt]
\psaxes[linewidth=1.25pt,labelFontSize=\scriptstyle,Dy=5]{->}(0,0)(-5,-4)(1.5,22)
\psplot[plotpoints=2000,linewidth=1.25pt,linecolor=red]{-5}{1}{10 x mul 5 add 2.71828 x exp mul}
\psplot[plotpoints=1000,linewidth=1.25pt,linestyle=dotted]{-0.7}{1}{15 x mul 5 add}
\psplot[plotpoints=1000,linewidth=1.25pt,linestyle=dotted]{-5}{0}{0.82085 x  mul 3.694 add neg}
\psdots(1,20)(-2.5,-1.6417)(0,5)
\uput[r](0,5){A}\uput[r](1,20){B}\uput[d](-2.5,-1.6417){C}\uput[l](0.6,19){$\red \mathcal{C}_f$}
\end{pspicture*}}
\medskip

\begin{enumerate}
\item On peut affirmer que:
	\begin{enumerate}
		\item $f'(-0,5) = 0$
		\item si $x \in ]- \infty~;~-0,5[$, alors $f'(x) < 0$
		\item $f'(0) = 15$
		\item la fonction dérivée $f'$ ne change pas de
signe sur $\R$.
	\end{enumerate}
\item On admet que la fonction $f$ représentée ci-dessus est définie sur $\R$ par 

$f(x) = (ax + b)\text{e}^x$, où $a$ et $b$ sont deux nombres réels et que sa courbe coupe l'axe des abscisses en son point de coordonnées $(-0,5~;~ 0)$.

On peut affirmer que:
	\begin{enumerate}
		\item $a = 10$ et $b = 5$
		\item $a = 2,5$ et $b = -0,5$
		\item $a = -1,5$ et $b = 5$
		\item $a = 0$ et $b = 5$
	\end{enumerate}
\item On admet que la dérivée seconde de la fonction $f$ est définie sur $\R$ par : 

\[f''(x) = (10x + 25)\text{e}^x.\]

On peut affirmer que :
	\begin{enumerate}
		\item La fonction $f$ est convexe sur $\R$
		\item La fonction $f$ est concave sur $\R$
		\item Le point C est l'unique point d'inflexion de $\mathcal{C}_f$
		\item $\mathcal{C}_f$ n'admet pas de point d'inflexion
	\end{enumerate}
\item On considère deux suites $\left(U_n\right)$ et  $\left(V_n\right)$ définies sur $\N$ telles que : 

\setlength\parindent{1cm}
\begin{itemize}
\item[$\bullet~~$] pour tout entier naturel $n$,\, $U_n \leqslant V_n$;
\item[$\bullet~~$] $\displaystyle\lim_{n \to+ \infty}  V_n= 2$.
\end{itemize}
\setlength\parindent{0cm}

On peut affirmer que:

	\begin{enumerate}
		\item la suite $\left(U_n\right)$ converge 
		\item pour tout entier naturel $n$,\, $V_n \leqslant 2$		
		\item la suite $\left(U_n\right)$ diverge
		\item la suite $\left(U_n\right)$ est majorée
	\end{enumerate} 
\end{enumerate}

\bigskip

\end{document}
